% !TEX root = Appendix_3.tex
% =========================================================================
% APPENDIX 3 -- FORMAL PROOFS AND LEAN VERIFICATION
% Econometrica / Econometric Society supplement format.
% Formal proofs supporting the identification argument. Replication: Appendix_3/
% =========================================================================
\documentclass[ecta,nameyear,final,supplement]{econsocart}

\usepackage{bm}
\usepackage{booktabs}
\usepackage{array}
\RequirePackage[colorlinks,citecolor=blue,linkcolor=blue,urlcolor=blue,pagebackref]{hyperref}
\urlstyle{same}

\startlocaldefs
\numberwithin{equation}{section}
\numberwithin{table}{section}

\newtheorem{theorem}{Theorem}[section]
\newtheorem{proposition}[theorem]{Proposition}
\newtheorem{lemma}[theorem]{Lemma}
\newtheorem{corollary}[theorem]{Corollary}
\theoremstyle{definition}
\newtheorem{definition}[theorem]{Definition}
\newtheorem{remark}[theorem]{Remark}

\newcommand{\E}{\mathbb{E}}
\newcommand{\ProbE}{\mathbb{P}_{\mathrm{E}}}
\newcommand{\ProbI}[1]{\mathbb{P}_{\mathrm{I},#1}}
\newcommand{\ProbC}[1]{\mathbb{P}_{\mathrm{C},#1}}
\newcommand{\Unif}{\mathrm{Unif}}

\newenvironment{fnnote}{%
  \begin{minipage}[t]{\linewidth}\sloppy
}{\end{minipage}}

\makeatletter
\@ifpackageloaded{etoolbox}{\AtEndEnvironment{definition}{\par\nobreak{\hfill\ensuremath{\square}\par}}}{}
\makeatother

\renewcommand{\thesection}{S\arabic{section}}
% Compact AMS-style labels: S{sec}.{sub}.{subsub}.{para}; 0 if level unused
\makeatletter
\newcommand{\@appzero}[1]{\ifnum\value{#1}=0 0\else\arabic{#1}\fi}
\newcommand{\@appsec}{S\arabic{section}}
\newcommand{\@appsub}{\@appzero{subsection}}
\newcommand{\@appsubsub}{\@appzero{subsubsection}}
\renewcommand{\thesubsection}{\@appsec.\@appsub.0.0}
\renewcommand{\thesubsubsection}{\@appsec.\@appsub.\@appsubsub.0}
\renewcommand{\theparagraph}{\@appsec.\@appsub.\@appsubsub.\arabic{paragraph}}
\setcounter{secnumdepth}{4}
\@addtoreset{subsection}{section}
\@addtoreset{subsubsection}{subsection}
\@addtoreset{paragraph}{section}
\makeatother

% APA 7-style table notes (Publication Manual, \S7.21)
\newcommand{\tabnote}{\par\addvspace{0.4em}\noindent\emph{Note.}}
\newcommand{\tabsource}{\ \textit{Source:} Analytical proofs in this supplement;
  Lean identifiers in \path{lean/Appendix3Proofs.lean}. Software: Lean~4
  (\path{lean-toolchain}); Mathlib via Lake.}
\endlocaldefs

\begin{document}

\begin{frontmatter}

\title{Supplement to ``Identifying Rational Types in Unknown Environments''\\
Appendix 3: Formal Proofs and Lean Verification}
\runtitle{Formal Proofs Supplement}

\begin{aug}
\author[add1,add2]{\fnms{Humberto}~\snm{Bernal}}
\address[add1]{%
\orgdiv{Economics Program},
\orgname{Universidad Colegio Mayor de Cundinamarca}}
\address[add2]{%
\orgdiv{Ph.D. in Economics},
\orgname{Universidad de los Andes}}
\ead[label=e1]{hbernal@uniandes.edu.co}
\ead[label=e2]{hbernalc@universidadmayor.edu.co}
\end{aug}

\begin{funding}
Economist, professor of the Economics Program at the Universidad Colegio
Mayor de Cundinamarca and graduate of the Ph.D.\ in Economics from the
Universidad de los Andes, Colombia.
\end{funding}

\end{frontmatter}

\noindent
This supplement presents in full detail the formal statements and proofs
supporting the identification argument in the main paper. It is organized in
eight sections that mirror the logical structure of the model: stochastic
trends (Section~\ref{sec:main}), public information and Bayesian learning
(Section~\ref{sec:info-bayes}), the stochastic implementation mechanism
(Section~\ref{sec:mdg}), asymptotic type identification
(Section~\ref{sec:id-asintotica}), competitive risks and operational pressure
(Section~\ref{sec:riesgos}), policy implications (Section~\ref{sec:politica}),
equilibrium implementability (Section~\ref{sec:equilibrio}), and the abstract
Kakutani identification criterion (Section~\ref{sec:id-abstracta}), plus
machine-checked verification (Section~\ref{sec:lean-formalization}). Analytical
proofs are self-contained in this document; algebraic steps verified in Lean~4
are reproduced from the \texttt{Appendix\_3/} replication package
(Section~\ref{sec:repro}). Classical measure-theoretic inputs (Kakutani,
Blackwell--Dubins, fixed-point existence) remain proved here and are not
re-derived inside the proof assistant.

% ============================================================================
\section{Results from the main paper}
\label{sec:main}
% ============================================================================

\noindent The following result is imported from
\textit{Identifying Rational Types in Unknown Environments: Evidence from
Kidnapping in Colombia}. It operates in the time-series domain and is
independent of the game-theoretic blocks below. Lemma~\ref{lem:I0-closed}
supplies the integration-order step.

% ---------------------------------------------------------------------------

\begin{lemma}[$I(0)$ closed under linear combinations]
\label{lem:I0-closed}
If $X_t\sim I(0)$ and $Y_t\sim I(0)$, then $aX_t+bY_t\sim I(0)$ for all
$a,b\in\mathbb{R}$.
\end{lemma}

\begin{proof}
Covariance stationarity is preserved under finite linear combinations: for all
$s,t$ and all $a,b\in\mathbb{R}$,
\[
\E\bigl[(aX_{t+s}+bY_{t+s})-(aX_t+bY_t)\bigr]=0,
\]
and the autocovariance function of $aX_t+bY_t$ depends on lags only through
$t-s$. Hence $aX_t+bY_t$ is integrated of order zero.
\end{proof}

\bigskip

\begin{corollary}[Equivalence of stochastic trends]
\label{cor:trend_equivalence}
Let $\{y_t\}$ be a time series integrated of order one, $I(1)$. Suppose the
existence of two valid decompositions:
\[
y_t = P_t^{(1)} + u_t^{(1)} = P_t^{(2)} + u_t^{(2)},
\]
where $P_t^{(1)}$ and $P_t^{(2)}$ are $I(1)$ processes and $u_t^{(1)}$,
$u_t^{(2)}$ are stationary $I(0)$ processes. Then:
\[
P_t^{(1)} - P_t^{(2)} \sim I(0).
\]
Consequently, all valid stochastic trends derived from the same observed
series are equivalent up to a stationary component.
\end{corollary}

\begin{proof}
From the two decompositions,
$P_t^{(1)}-P_t^{(2)}=u_t^{(2)}-u_t^{(1)}$. Since $u_t^{(1)}$ and $u_t^{(2)}$
are $I(0)$, Lemma~\ref{lem:I0-closed} with $a=1$ and $b=-1$ yields
$P_t^{(1)}-P_t^{(2)}\sim I(0)$. Any pair of valid stochastic trends extracted
from the same $I(1)$ series therefore differs by a transitory $I(0)$ component.
\end{proof}

% ============================================================================
\section{Public information, likelihood, and Bayesian learning}
\label{sec:info-bayes}
% ============================================================================

\noindent The three definitions below construct the public filtration, the
likelihood kernel, and the belief process. The stability lemma closes the
block.

% ---------------------------------------------------------------------------

\begin{definition}[Public signal and filtration]
\label{def:zeta-filtration}
\noindent\textbf{(i) Public signal $\zeta_t$.}
The public signal $\zeta_t$ is the random vector that summarizes what is
observable at the close of period~$t$:
\[
\zeta_t=(\alpha_t,\gamma_t,\,\tilde{a}_t^F,\tilde{a}_t^K,\tilde{a}_t^S,\,m_t,d_t)
\in
\mathcal{Z}_t\subseteq\mathbb{R}^2\times\mathcal{A}^F\times\mathcal{A}^K
\times\mathcal{A}^{S,\mathrm{disc}}\times\mathcal{M}_{\mathrm{out}}\times\{0,1\}.
\]

\noindent\textbf{(ii) Public filtration $\{\mathcal{F}_t^{\mathrm{pub}}\}$.}
Set $\mathcal{F}_0^{\mathrm{pub}}:=\sigma(z,\theta_V)$. For $t\ge 1$,
\[
\mathcal{F}_t^{\mathrm{pub}}:=\sigma(z,\zeta_0,\ldots,\zeta_{t-1}),
\]
that is, the $\sigma$-algebra generated by the initial location and by the
signals observed strictly before period~$t$ begins.

\noindent\textbf{(iii) Measurability.}
Each $\zeta_t:(\Omega,\mathcal{F})\to(\mathcal{Z}_t,\mathcal{B}(\mathcal{Z}_t))$
is measurable, and $\sigma(z,\zeta_0,\ldots,\zeta_{t-1})$ is well defined as a
sub-$\sigma$-algebra of $\mathcal{F}$.
\end{definition}

\bigskip

\begin{definition}[Predictive likelihood under a counterfactual type]
\label{def:predictive-likelihood}
Fixing a PBE with profile $(\chi,s^F,s^K)$ and a counterfactual type
$\theta\in\Theta_K$, the \textbf{predictive likelihood} of period~$t$ is
\begin{equation}
\label{eq:likelihood-zeta}
\ell_t(\theta;J)
:=\ProbC{\theta}\bigl(\zeta_t\in J\mid\mathcal{F}_t^{\mathrm{pub}}\bigr),
\qquad J\in\mathcal{B}(\mathcal{Z}_t).
\end{equation}
The mapping $J\mapsto\ell_t(\theta;J)$ is a \textbf{probability kernel} that
carries the information state $\mathcal{F}_t^{\mathrm{pub}}$ to the law of
$\zeta_t$ under type~$\theta$.
\end{definition}

\bigskip

\begin{definition}[Beliefs and PBE-consistent Bayesian updating]
\label{def:beliefs-martingale}
Let $\mu_0\in\Delta(\Theta_K)$ be the \textbf{common prior} with full support,
$\mu_0(\theta\mid z,\theta_V)>0$ for all $\theta\in\Theta_K$.
For $t\ge 1$, let $\mu_t\in\Delta(\Theta_K)$ be the distribution obtained by
iterated Bayes from $\mu_0$, using the realizations of
$(\zeta_0,\ldots,\zeta_{t-1})$ and the predictive kernels
$\ell_s(\theta';J)$ induced by the same PBE.

The process $\{\mu_t\}_{t\ge 0}$ is \textbf{consistent with the PBE} if, under
the ex-ante law $\ProbE$,
\begin{equation}
\label{eq:beliefs-ebp-consistency}
\mu_t(\theta)
=\ProbE\bigl(\theta_K=\theta\mid\mathcal{F}_t^{\mathrm{pub}}\bigr)
\qquad\forall\,\theta\in\Theta_K,\;\forall\,t\ge 1,
\end{equation}
interpreted $\ProbE$-almost surely upon fixing regular versions of the
conditional probabilities.
\end{definition}

\bigskip

\begin{lemma}[Asymptotic stability of beliefs. Convergence theorem of
\citeauthor{Doob1953Stochastic}, \citeyear{Doob1953Stochastic}]
\label{lem:asymptotic-beliefs}
Under the ex-ante law $\ProbE$ of
Definition~\ref{def:beliefs-martingale}, the process
$\{\mu_t(\theta)\}_{t\ge 0}$ is a \textbf{bounded martingale} that converges
$\ProbE$-almost surely to a limit random variable
$Z_\theta\in[0,1]$ $\ProbE$-a.s.
\end{lemma}

\begin{proof}
The proof verifies the conditions of the convergence theorem of
\citet{Doob1953Stochastic,Williams1991Martingales,Shiryaev2019Probability}.

\smallskip
\noindent\textbf{1.~Martingale property.}
Let $M=\mathbf{1}\{\theta_K=\theta\}$. By equilibrium consistency
(Definition~\ref{def:beliefs-martingale}),
\[
\mu_t(\theta)=\E^{\ProbE}[M\mid\mathcal{F}_t^{\mathrm{pub}}]\quad\ProbE\text{-a.s.}
\]
As $\mathcal{F}_t^{\mathrm{pub}}\subseteq\mathcal{F}_{t+1}^{\mathrm{pub}}$,
the \textbf{law of iterated expectations} yields
\[
\E^{\ProbE}[\mu_{t+1}(\theta)\mid\mathcal{F}_t^{\mathrm{pub}}]
=\E^{\ProbE}\!\bigl[\E^{\ProbE}[M\mid\mathcal{F}_{t+1}^{\mathrm{pub}}]
                     \mid\mathcal{F}_t^{\mathrm{pub}}\bigr]
=\E^{\ProbE}[M\mid\mathcal{F}_t^{\mathrm{pub}}]
=\mu_t(\theta)\quad\ProbE\text{-a.s.}
\]
The process is a martingale under $\ProbE$.

\smallskip
\noindent\textbf{2.~Pointwise and uniform $L^1$ bound.}
As $\mu_t(\theta)$ is a probability mass,
$0\le\mu_t(\theta)\le 1$ for all $t\ge 0$ (outside a
$\ProbE$-null set). In particular,
$\sup_{t\ge 0}\E^{\ProbE}[|\mu_t(\theta)|]\le 1$, which provides the
uniform $L^1$ bound required by Doob's theorem.

\smallskip
\noindent\textbf{3.~Almost sure convergence.}
Once the martingale property and the uniform $L^1$ bound are verified,
\textbf{Doob's convergence theorem} guarantees the existence of
$Z_\theta$ such that
\[
\lim_{t\to\infty}\mu_t(\theta)=Z_\theta\quad\ProbE\text{-a.s.}
\]
The pointwise bound $0\le\mu_t(\theta)\le 1$ implies $Z_\theta\in[0,1]$
$\ProbE$-a.s.
\end{proof}

% ============================================================================
\section{Stochastic implementation mechanism (MDG)}
\label{sec:mdg}
% ============================================================================

\noindent This section establishes the fundamental definitions of the
stochastic implementation layer (MDG), which transforms strategic intentions
into materialized actions, and of the law governing the physical outcome that
determines the observable results of the episode.

% ---------------------------------------------------------------------------

\begin{definition}[MDG stochastic implementation filter]
\label{def:mano-dios-filtro}
Let $\mathcal{A}^i$ be the finite set of admissible actions of agent~$i$.
Given the plan $a_t^{i\ast}\in\mathcal{A}^i$ and the state vector $X_t$,
execution is determined by the \textbf{implementation law}
\[
\tilde{a}_t^i \sim \ProbI{i}\bigl(\cdot \mid a_t^{i\ast},\,X_t\bigr).
\]
The stochastic structure is formalized on a measurable space
$(\Omega,\mathcal{F})$. The measurable function
$\mathcal{H}_t^i:\mathcal{A}^i\times\Omega\to\mathcal{A}^i$ materializes the
sequential draw: $\tilde{a}_t^i(\omega)=\mathcal{H}_t^i(a_t^{i\ast},\omega)$.
The \textbf{aggregate profile of executed actions} is
\begin{equation}
\label{eq:tilde-A-t}
\tilde{A}_t = (\tilde{a}_t^K,\,\tilde{a}_t^S,\,\tilde{a}_t^F).
\end{equation}
The following regularity assumptions are imposed.
\begin{enumerate}
  \item[\textup{(i)}] \textbf{Plan dominance and error support.}
  For all $t<\tau$, $\ProbI{i}$ concentrates most of its mass on
  $a_t^{i\ast}$ and assigns strictly positive probability to the remaining
  actions. Under the \textbf{full support} hypothesis,
  $\ProbI{i}(\tilde{a}_t^i=a\mid a_t^{i\ast},X_t)>0$ for all
  $a\in\mathcal{A}^i$.
  \item[\textup{(ii)}] \textbf{Institutional convergence.}
  In the long-horizon regime one may impose
  \[
  \lim_{t\to\infty}\ProbE\bigl(\tilde{a}_t^i=a_t^{i\ast}\mid
    \mathcal{F}_t^{\mathrm{pub}}\bigr)=1,\qquad\forall\,i,
  \]
  which is satisfied when the posterior entropy $H(\mu_t)\to 0$.
\end{enumerate}
The implementation law adopts the \textbf{hybrid-temperature logit}
structure:
\begin{equation}
\label{eq:logit-hybrid}
\ProbI{i}\bigl(\tilde{a}_t^i=a\mid a_t^{i\ast},X_t\bigr)
=\frac{\exp\bigl(\mathbf{1}\{a=a_t^{i\ast}\}/T_t\bigr)}
      {\sum_{a'\in\mathcal{A}^i}\exp\bigl(\mathbf{1}\{a'=a_t^{i\ast}\}/T_t\bigr)},
\end{equation}
with hybrid temperature
\begin{equation}
\label{eq:temperatura-piso}
T_t = T_0\!\left[\frac{H(\mu_t)}{H(\mu_0)}\right]
      \max\!\bigl\{e^{-\eta_{\mathrm{cal}}t},\,\underline{c}\bigr\},
\end{equation}
where $T_0>0$ is the initial temperature and $\eta_{\mathrm{cal}}>0$ is the
calendar decay rate. Shannon entropy is
$H(\mu_t)=-\sum_\theta\mu_t(\theta)\ln\mu_t(\theta)$, and
$\underline{c}\in(0,1]$ is the floor of structural noise.
\end{definition}

\begin{remark}[Axiomatic justification of the logit as the unique implementation law]
\label{rem:axiomas-logit}
Equation~\eqref{eq:logit-hybrid} is the \emph{unique} distribution over
$\mathcal{A}^i$ that simultaneously satisfies the three axioms:
\begin{enumerate}
  \item[\textup{(A1)}] \textbf{Full support:}
  $\ProbI{i}(a\mid a_t^{i\ast},X_t)>0$ for all $a\in\mathcal{A}^i$.
  \item[\textup{(A2)}] \textbf{Dominance:}
  $\ProbI{i}(a_t^{i\ast}\mid a_t^{i\ast},X_t)>\ProbI{i}(a\mid a_t^{i\ast},X_t)$
  for all $a\neq a_t^{i\ast}$.
  \item[\textup{(A3)}] \textbf{Maximum entropy:}
  Among all distributions satisfying~(A1) and~(A2),
  $\ProbI{i}$ maximizes the Shannon entropy
  $H=-\sum_a p_a\ln p_a$ subject to the optimal action receiving fixed
  probability $1-\varepsilon_t\in(0,1)$.
\end{enumerate}
By the characterization theorem of \citet{Jaynes1957Information}, the unique
maximizer is the logit in~\eqref{eq:logit-hybrid}. The MDG process is thus
axiomatically justified as the unique implementation process compatible with
bounded rationality and maximum entropy.
\end{remark}

\bigskip

\begin{definition}[Law of the physical outcome]
\label{def:desenlace-fisico-mt}
The \textbf{physical outcome} $m_t$ is an observable categorical variable on
the support $\mathcal{M}_{\mathrm{out}}$, where
$\mathcal{M}_{\mathrm{cont}}=\{\mathrm{cont}\}$,
$\mathcal{M}_{\mathrm{term}}=\{\mathrm{pay},\mathrm{res},\mathrm{rel},\mathrm{kill}\}$
and $\mathcal{M}_{\mathrm{out}}=\mathcal{M}_{\mathrm{cont}}\cup\mathcal{M}_{\mathrm{term}}$.
The label $\mathrm{cont}$ represents the continuation of captivity and
$j\in\{1,2,3,4\}$ closure by competitive cause~$j$. For $j\in\{1,2,3\}$, the
effective intensity $\tilde\lambda_j(t\mid\theta_K)$ carries the type-specific
parameter $\beta_{K,j}(\theta_K)$ in the Cox specification; cause $j=4$ enters
only through the baseline hazard (there is no $\beta_{K,4}$). Conditional on the
captor's type and the material state of the period
\[
\mathcal{C}_t(\theta_K)
=\bigl(\theta_K,\theta_F,\theta_V,\theta_S,z,
\tilde{a}_t^F,\tilde{a}_t^K,\tilde{a}_t^S,
\alpha_t^\ast,\gamma_t^\ast,p_{\mathrm{det},t}\bigr),
\]
the distribution of $m_t$ is defined from the effective intensities
$\tilde\lambda_j(t\mid\theta_K)$. Denote by
\[
Q_t(m\mid\theta_K,\mathcal C_t,\alpha_t^\ast,\gamma_t^\ast)
:=\ProbE(m_t=m\mid\theta_K,\mathcal C_t,\alpha_t^\ast,\gamma_t^\ast)
\]
the transition kernel of the physical outcome, and by $\mathsf{G}_t$ the
measurable rule that materializes a concrete realization of $m_t$ from that
kernel and a uniform draw. The \textbf{continuation probability} is
\[
Q_t(\mathrm{cont}\mid\theta_K,\mathcal{C}_t,\alpha_t^\ast,\gamma_t^\ast)
=\exp\!\left(-\sum_{j=1}^{4}\tilde\lambda_j(t\mid\theta_K)\,\Delta t\right).
\]
The daily closure total and the relative weight of each cause are
\[
q_t(\theta_K)=1-p_{\mathrm{Cont},t},\quad
\xi_j=\frac{\tilde\lambda_j}{\sum_\ell\tilde\lambda_\ell},\quad
\bar h_j(t\mid\theta_K)=q_t\,\xi_j.
\]
The materialization of $m_t$ is obtained by \textbf{inverse transform}:
let $v_t\sim\Unif(0,1)$; then
\[
m_t=r
\iff
\sum_{\ell<r}Q_t(\ell\mid\theta_K,\mathcal{C}_t,\alpha_t^\ast,\gamma_t^\ast)\le v_t
<\sum_{\ell\le r}Q_t(\ell\mid\theta_K,\mathcal{C}_t,\alpha_t^\ast,\gamma_t^\ast),
\]
and the likelihood of $m_t$ in the Bayesian update is
\[
\ProbE(m_t\mid\theta_K,\mathcal{C}_t)
=p_{\mathrm{Cont}}^{\mathbf{1}\{m_t=\mathrm{cont}\}}
\prod_{j=1}^{4}\bar h_j^{\mathbf{1}\{m_t=j\}}.
\]
\end{definition}

% ============================================================================
\section{Asymptotic type identification}
\label{sec:id-asintotica}
% ============================================================================

\noindent This section defines the prolonged captivity event and the notion of
asymptotic identification, states the structural identifiability lemma (bridge
between the parameters $\beta_{K,j}$ and the singularity of laws), and
concludes with the posterior consistency theorem.

\medskip
\noindent\textbf{Notation.}
Throughout Sections~\ref{sec:id-asintotica} and~\ref{sec:id-abstracta},
$\theta_K^\ast$ denotes the realized captor type. When conditioning on the law
of the true type, the notation $\theta^\ast:=\theta_K^\ast$ and
$\mu_C:=\ProbC{\theta^\ast}(\cdot\mid\mathcal{T})$ is adopted.

% ---------------------------------------------------------------------------

\begin{definition}[Prolonged captivity event]
\label{def:event-T-long}
Let $\tau$ be the stopping time of the episode. The
\textbf{prolonged captivity event} is
\[
\mathcal{T}
:=\{\tau=\infty\}
=\bigl\{\omega\in\Omega:m_t(\omega)\in\mathcal{M}_{\mathrm{cont}}\;\forall\,t\ge 0\bigr\}.
\]
In $\mathcal{T}$ the episode admits no terminal outcome in any finite period
and generates an infinite sequence of public signals
$\{\zeta_t\}_{t\ge 0}$.
\end{definition}

\bigskip

\begin{definition}[Asymptotic identification in the event $\mathcal{T}$]
\label{def:asymptotic-id}
The model is \textbf{asymptotically identifiable} in $\mathcal{T}$
(Definition~\ref{def:event-T-long}) if
$\ProbC{\theta_K^\ast}(\mathcal{T})>0$ and, for every $\theta\in\Theta_K$ with
$\theta\neq\theta_K^\ast$, at least one of the following conditions holds.
\begin{enumerate}
  \item[\textup{(a)}] \textbf{Exclusion by survival.}
  $\ProbC{\theta}(\mathcal{T})=0$.
  \item[\textup{(b)}] \textbf{Tail singularity of signals.}
  $\ProbC{\theta}(\mathcal{T})>0$ and the laws of $\{\zeta_t\}_{t\ge 0}$
  restricted to $\mathcal{T}$ under $\ProbC{\theta}$ and under
  $\ProbC{\theta_K^\ast}$ are \textbf{mutually singular}.
\end{enumerate}
\end{definition}

\begin{remark}[Role of Lemma~\ref{lem:hazard-separation} in the analysis]
Theorem~\ref{thm:posterior-consistency} requires asymptotic identification
(Definition~\ref{def:asymptotic-id}), a condition formulated in terms of
laws of entire processes. Lemma~\ref{lem:hazard-separation} is the
\textbf{operational bridge} between that abstract hypothesis and the
parameters $\beta_{K,j}$ already introduced in the mechanism: it allows one
to verify in the practice of the model when two distinct types generate
asymptotically distinguishable trajectories. Without this result, the learning
section would rest on a high-level condition difficult to connect with the
proportional parametrization of the $\lambda_j$.
\end{remark}

\bigskip

\begin{lemma}[Structural identifiability and asymptotic learning]
\label{lem:hazard-separation}
Suppose $|\Theta_K|<\infty$. Let $Y_t:=(m_t,d_t)$ be the daily public record.
Under a PBE, $p_\theta(\cdot\mid\mathcal{F}_t^{\mathrm{pub}})$ denotes the
mass function of $Y_t$ conditioned on public information at the opening of~$t$
under type~$\theta$.

\smallskip
\noindent\textbf{Part~(i) --- local sufficiency
\citep{Cox1972,Bradburn2003}.}
Let $\theta,\theta'\in\Theta_K$ be distinct. If there exists $j\in\{1,2,3\}$ such that
$\beta_{K,j}(\theta)\neq\beta_{K,j}(\theta')$ (causes $j\in\{1,2,3\}$ only; see
Definition~\ref{def:desenlace-fisico-mt}) and no exact compensations arise
between the masses of $Y_t$ (\emph{knife-edges}),
then there exists $\varepsilon=\varepsilon(\theta,\theta',t)>0$ such that
\begin{equation}
\label{eq:TV-separation}
\bigl\|p_\theta(\cdot\mid\mathcal{F}_t^{\mathrm{pub}})
-p_{\theta'}(\cdot\mid\mathcal{F}_{t}^{\mathrm{pub}})\bigr\|_{\mathrm{TV}}
\ge\varepsilon.
\end{equation}

\smallskip
\noindent\textbf{Part~(ii) --- asymptotic sufficiency
\citep{Kakutani1948,Shiryaev2019Probability}.}
Fixing a policy trajectory $(\chi_s)_{s\ge 0}$, suppose that
conditional on $\mathcal{T}$ the $\{Y_t\}_{t\ge 0}$ are independent. For
each pair $\theta\neq\theta'$, let $\mathcal{S}_{\theta,\theta'}\subseteq\mathbb{N}$
be the set of periods where~\eqref{eq:TV-separation} holds. If
$\mathcal{S}_{\theta,\theta_K^\ast}$ is infinite $\ProbC{\theta_K^\ast}$-a.s.\
and there exists $\underline\varepsilon_\theta>0$ such that
$\|p_\theta^{(t)}-p_{\theta_K^\ast}^{(t)}\|_{\mathrm{TV}}
\ge\underline\varepsilon_\theta$ for all $t\in\mathcal{S}_{\theta,\theta_K^\ast}$,
then the laws of $(Y_t)_{t\ge 0}$ conditioned on $\mathcal{T}$ under
$\ProbC{\theta}$ and under $\ProbC{\theta_K^\ast}$ are
\textbf{mutually singular}.
\end{lemma}

\begin{proof}
The proof identifies the causal mechanism within the day, translates it to a
statistical bound in total variation, and couples the repeated experiment to
the criterion of \citet{Kakutani1948}.

\smallskip
\noindent\textbf{1. Same control menu, distinct criminal technology.}
Fix $t$ and $\mathcal{F}_t^{\mathrm{pub}}$. In the counterfactual that changes
only the type ($\theta$ vs.\ $\theta'$), keeping equal the executed triple
$\tilde{A}_t$, the instruments $(\alpha_t^\ast,\gamma_t^\ast)$, and the
maturation filter $M(t)$, the proportional parametrization implies
\[
\frac{\lambda_j(t\mid\theta)}{\lambda_j(t\mid\theta')}
=\exp\!\bigl(\beta_{K,j}(\theta)-\beta_{K,j}(\theta')\bigr)\neq 1.
\]

\smallskip
\noindent\textbf{2. From daily intensity to distinct predictions over $Y_t$.}
The proportional difference in $\lambda_j(t)$ is projected onto the aggregate
exponential of $q(t)$, the causal mix $\xi_j(t)$, or the marginal hazards
$\bar h_j(t)=q(t)\xi_j(t)$. Except for \emph{knife-edges}, the marginal of $m_t$
cannot coincide under $\theta$ and $\theta'$, and the two mass functions differ
in total variation, verifying~\eqref{eq:TV-separation}.

\smallskip
\noindent\textbf{3. Repetition with uniform margin and Kakutani criterion.}
Conditional on the fixed policy and on $\mathcal{T}$, the law of the sequence is
a product of daily experiments
$\bigotimes_{t\ge 0}p_\theta^{(t)}$ versus
$\bigotimes_{t\ge 0}p_{\theta_K^\ast}^{(t)}$. Part~(ii) then follows from
Lemma~\ref{lem:general-id-appendix} (Section~\ref{sec:id-abstracta}): recurrent
uniform total-variation separation on an infinite set forces the Bhattacharyya
product affinity to zero and, by the criterion of \citet{Kakutani1948}, mutual
singularity of the infinite product measures
\citep{Shiryaev2019Probability}.
\end{proof}

\bigskip

\begin{theorem}[Identification and consistency of true types]
\label{thm:posterior-consistency}
Consider a prior $\mu_0\in\Delta(\Theta_K)$ with full support and a
PBE-consistent Bayesian updating
(Definition~\ref{def:beliefs-martingale}). Let $\theta^\ast:=\theta_K^\ast$ be the
realized type and $\mu_C:=\ProbC{\theta^\ast}(\cdot\mid\mathcal{T})$ the law
of the true type restricted to the event $\mathcal{T}$
(Definition~\ref{def:event-T-long}).

If the mechanism is \textbf{asymptotically identifiable} in $\mathcal{T}$
(Definition~\ref{def:asymptotic-id}), the following properties are satisfied
$\mu_C$-almost surely:
\begin{enumerate}
  \item \textbf{Exclusion of false hypotheses:}
  $\mu_t(\theta)\xrightarrow{t\to\infty}0$ for all $\theta\neq\theta^\ast$.
  \item \textbf{Concentration on the true type:}
  $\mu_t(\theta^\ast)\xrightarrow{t\to\infty}1$.
\end{enumerate}
Consequently, in $\mathcal{T}$ the sequence $\{\mu_t\}$ converges weakly
to the Dirac measure $\delta_{\theta^\ast}$.
\end{theorem}

\begin{proof}
The proof separates the general probabilistic argument from the structural
identification assumption.

\paragraph{1. Stability of beliefs.}
By equilibrium consistency,
$\mu_t(\theta)=\E^{\ProbE}[\mathbf{1}\{\theta_K=\theta\}\mid
\mathcal{F}_t^{\mathrm{pub}}]$ $\ProbE$-a.s. Lemma
~\ref{lem:asymptotic-beliefs} guarantees that $\{\mu_t(\theta)\}$ is a
bounded martingale with limit $Z_\theta\in[0,1]$ $\ProbE$-a.s. The
transfer to $\mu_C$-a.s.\ is justified by disintegration: the mixture
$\ProbE=\sum_\vartheta\mu_0(\vartheta)\ProbC{\vartheta}$ implies that every
$\ProbE$-null set is $\ProbC{\theta^\ast}$-null (since
$\mu_0(\theta^\ast)>0$), and domination is preserved upon restricting to
$\mathcal{T}$.

\paragraph{2. Identification and exclusion of false types.}
Fix $\theta\neq\theta^\ast$.

\emph{Case} $\ProbC{\theta}(\mathcal{T})=0$: the trajectory under $\mu_C$
belongs to an event that the false type cannot generate; its posterior
likelihood is null in the conditioned regime.

\emph{Case} $\ProbC{\theta}(\mathcal{T})>0$: the singularity of
Definition~\ref{def:asymptotic-id}(b) implies that there exists a tail event
$A_\theta$ with
$\ProbC{\theta^\ast}(A_\theta\mid\mathcal{T})=1$ and
$\ProbC{\theta}(A_\theta\mid\mathcal{T})=0$.
Let $R_t^\theta$ be the likelihood ratio of the false type against the true
type on finite horizons. Tail singularity forces
$R_t^\theta\to 0$ $\mu_C$-a.s.\
\citep{BlackwellDubins1962,Schervish1995Theory}.
By Bayes' formula with a full-support prior,
\[
\mu_t(\theta)
=\frac{\mu_0(\theta)\,R_t^\theta}
      {\mu_0(\theta^\ast)+\sum_{\theta'\neq\theta^\ast}\mu_0(\theta')R_t^{\theta'}}
\le\frac{\mu_0(\theta)}{\mu_0(\theta^\ast)}\,R_t^\theta,
\]
and since $\mu_0(\theta^\ast)>0$, it follows that
$\mu_t(\theta)\to 0$ $\mu_C$-a.s.\ for all $\theta\neq\theta^\ast$.

\paragraph{3. Closure by normalization.}
For all $t$, $\sum_{\theta\in\Theta_K}\mu_t(\theta)=1$. As $\Theta_K$ is
finite and all false coordinates converge to zero $\mu_C$-a.s.,
\[
1=\lim_{t\to\infty}\sum_\theta\mu_t(\theta)
=Z_{\theta^\ast}+\sum_{\theta\neq\theta^\ast}0,
\]
hence $Z_{\theta^\ast}=1$ $\mu_C$-a.s., that is,
$\mu_t(\theta^\ast)\to 1$ $\mu_C$-a.s. On a finite simplex, coordinate
convergence is equivalent to weak convergence $\mu_t\to\delta_{\theta^\ast}$.
\end{proof}

\begin{remark}[Economic reading and link with $\beta_{K,j}$]
\label{rem:lectura-beta}
Definition~\ref{def:asymptotic-id} requires that the design of the game
(policy + incentives + risk technology) produce long trajectories incompatible
with false types. Distinct shifts
$\beta_{K,j}(\theta)\neq\beta_{K,j}(\theta')$ favor the singularity of
the laws of $(m_t,d_t)$. If the equilibrium pools types, Lemma
~\ref{lem:asymptotic-beliefs} still holds but Theorem
~\ref{thm:posterior-consistency} does not: the limit may split mass among
observationally equivalent types.
\end{remark}

\begin{remark}[Finite horizon]
\label{rmk:horizonte-finito}
If $\tau<\infty$ almost surely, $\mathcal{T}$ has measure zero and the
statement of Theorem~\ref{thm:posterior-consistency} is vacuous. The relevant
conclusion is then only the martingale convergence of Lemma
~\ref{lem:asymptotic-beliefs}: beliefs stabilize at the finite limit of
observations.
\end{remark}

% ============================================================================
\section{Competitive risks and operational pressure}
\label{sec:riesgos}
% ============================================================================

\noindent The following proposition uses directly the proportional risk
structure, without requiring the equilibrium machinery of
Section~\ref{sec:equilibrio}. It depends only on the intensities
$\tilde\lambda_j$ introduced in Definition~\ref{def:desenlace-fisico-mt}.

% ---------------------------------------------------------------------------

\begin{proposition}[Expected duration of captivity and operational pressure]
\label{prop:expected-captivity-gamma}
Define the \textbf{net sensitivity} of type $\theta_K$ to optimal operational
pressure $\gamma_t^\ast$ as
\begin{equation}
\label{eq:kappa-c}
\kappa_h(\theta_K,t)
:=\bigl(\zeta_{\gamma,2}\,\tilde\lambda_2(t\mid\theta_K)
        +\zeta_{\gamma,3}\,\tilde\lambda_3(t\mid\theta_K)\bigr)
  -\zeta_{\gamma,1}\,\tilde\lambda_1(t\mid\theta_K).
\end{equation}
If $\kappa_h(\theta_K,t)>0$ for all $t$, the expected duration of captivity
$\mathbb{E}[\tau\mid\theta_K]$ is strictly decreasing in $\gamma_t^\ast$.
If $\kappa_h(\theta_K,t)<0$ for all $t$, it is strictly increasing. The model
generates the \textbf{testable differential prediction}:
\begin{equation}
\operatorname{sgn}\!\left(
  \frac{\partial\,\mathbb{E}[\tau\mid\theta_K]}{\partial\gamma_t^\ast}
\right)
=-\operatorname{sgn}\!\left(\kappa_h(\theta_K,t)\right),
\end{equation}
with the sign varying across types through the structural parameters
$\beta_{K,j}(\theta_K)$.
\end{proposition}

\begin{proof}
The proof proceeds in three steps: differentiation of daily survival,
propagation by the chain rule, and conclusion by dominated convergence.

\paragraph{Step 1. Duration identity and survival.}
By \textbf{Fubini's theorem} \citep{Fubini1907},
\begin{equation}
\label{eq:duracion-identidad}
\mathbb{E}[\tau\mid\theta_K]
=\sum_{t'=0}^{\infty}S(t'\mid\theta_K),
\end{equation}
where $S(t'\mid\theta_K):=\prod_{s=1}^{t'}p_{\mathrm{Cont},s\mid\theta_K}$,
with $S(0\mid\theta_K)=1$. Under the proportional parametrization,
\[
p_{\mathrm{Cont},s\mid\theta_K}
=\exp\!\Bigl(-\sum_{j=1}^{4}\tilde\lambda_j(s\mid\theta_K)\,\Delta t\Bigr),
\]
consistent with Definition~\ref{def:desenlace-fisico-mt}. Since only the causes
$j\in\{1,2,3\}$ respond to operational pressure (cause $j=4$ carries no
$\beta_{K,4}$ and does not depend on $\gamma_t^\ast$), differentiating with
respect to $\gamma_t^\ast$:
\begin{equation}
\label{eq:dpCont-dgamma}
\frac{\partial\,p_{\mathrm{Cont},t\mid\theta_K}}{\partial\gamma_t^\ast}
=-\kappa_h(\theta_K,t)\;p_{\mathrm{Cont},t\mid\theta_K}\;\Delta t.
\end{equation}

\paragraph{Step 2. Propagation to accumulated survival.}
For $t'\ge t$, only the factor $s=t$ in the product depends on $\gamma_t^\ast$;
by the chain rule:
\begin{equation}
\label{eq:dS-dgamma}
\frac{\partial S(t'\mid\theta_K)}{\partial\gamma_t^\ast}
=-\kappa_h(\theta_K,t)\;S(t'\mid\theta_K)\;\Delta t,\qquad t'\ge t.
\end{equation}
For $t'<t$, the derivative is zero.

\paragraph{Step 3. Expected duration.}
Differentiating~\eqref{eq:duracion-identidad} and applying
\eqref{eq:dS-dgamma} gives
\begin{equation}
\label{eq:dE-dgamma}
\frac{\partial\,\mathbb{E}[\tau\mid\theta_K]}{\partial\gamma_t^\ast}
=-\kappa_h(\theta_K,t)\,\Delta t\sum_{t'=t}^{\infty}S(t'\mid\theta_K).
\end{equation}
The interchange of sum and derivative is justified by the \textbf{dominated
convergence theorem}: the bound
$|\partial S(t')/\partial\gamma_t^\ast|\le|\kappa_h|\,S(t')\,\Delta t$ is
summable because $\mathbb{E}[\tau\mid\theta_K]<\infty$, guaranteed by the
uniform positivity of the baseline hazards.

As $S(t'\mid\theta_K)\ge 0$ and the tail $\sum_{t'\ge t}S(t')>0$:
\[
\operatorname{sgn}\!\left(
  \frac{\partial\,\mathbb{E}[\tau\mid\theta_K]}{\partial\gamma_t^\ast}
\right)
=-\operatorname{sgn}\!\left(\kappa_h(\theta_K,t)\right).\qedhere
\]
\end{proof}

% ============================================================================
\section{Policy implications}
\label{sec:politica}
% ============================================================================

\noindent This proposition requires the belief machinery of
Section~\ref{sec:info-bayes} and the equilibrium context of
Section~\ref{sec:equilibrio}: optimal pressure is calibrated in a Bayesian
manner.

% ---------------------------------------------------------------------------

\begin{proposition}[Optimal operational pressure and Bayesian type identification]
\label{prop:optimal-gamma-beliefs}
Suppose $\Theta_K=\{\theta_{low},\theta_{high}\}$ and fix a discrete branch
$b\in\{R,N\}$ in a neighborhood of the posterior belief under consideration. Let
$x_t=(\alpha_t,\gamma_t)$ and define the complete branch objective
\begin{equation}
\label{eq:obj-rama-completo}
J_t^b(x_t,\mu_t)
:=
\sum_{\theta\in\Theta_K}\mu_t(\theta)
L_t(b,\alpha_t,\gamma_t,\theta,\iota_t)
+\Pi_{t,b}^{S}(\alpha_t,\gamma_t;\mu_t)
-\psi_H\Delta H_t(\alpha_t,\gamma_t).
\end{equation}
If $(\alpha_t^\ast,\gamma_t^\ast)$ is an interior optimum of
$J_t^b(\cdot,\mu_t)$, the constraints $IC^K$, $IR^K$, and $IR^F$ do not locally
change regime, the Hessian
$H_t^b:=\nabla_{xx}^2J_t^b(x_t^\ast,\mu_t)$ is positive definite, and the net
marginal effect of increasing $\mu_t(\theta_{high})$ on the first-order
condition for $\gamma_t$ favors greater pressure, that is,
\begin{equation}
\label{eq:monotonia-neta-gamma}
e_\gamma'
\bigl(H_t^b\bigr)^{-1}
\nabla_{x\mu}^2J_t^b(x_t^\ast,\mu_t)
<0,
\qquad
e_\gamma=(0,1)',
\end{equation}
then the optimal operational pressure of that branch is locally increasing in
the posterior belief on the high-risk type:
\[
\frac{d\gamma_t^\ast}{d\mu_t(\theta_{high})}>0.
\]
The structural separation $\beta_{K,j}(\theta_{high})-\beta_{K,j}(\theta_{low})$
for $j\in\{2,3\}$ contributes to this condition through the ratio
\[
\frac{\tilde\lambda_j(t\mid\theta_{high})}{\tilde\lambda_j(t\mid\theta_{low})}
=\exp\!\bigl(\beta_{K,j}(\theta_{high})-\beta_{K,j}(\theta_{low})\bigr),
\quad j\in\{2,3\}.
\]
Monotonicity, however, depends on the joint net effect of material losses,
deviation penalties, and informational gains.
\end{proposition}

\begin{proof}
Write $\mu:=\mu_t(\theta_{high})$ and
$\mu_t(\theta_{low})=1-\mu$. Within the fixed discrete branch $b$, the interior
optimum $x_t^\ast(\mu)=(\alpha_t^\ast(\mu),\gamma_t^\ast(\mu))$ satisfies
\begin{equation}
\label{eq:foc-rama-completa}
\nabla_x J_t^b(x_t^\ast(\mu),\mu)=0.
\end{equation}
The hypothesis $H_t^b=\nabla_{xx}^2J_t^b(x_t^\ast,\mu)\succ0$ implies that the
local optimum is isolated and that the Jacobian of the first-order condition
with respect to $x_t$ is invertible. By the implicit function theorem,
$x_t^\ast(\mu)$ is locally differentiable and
\begin{equation}
\label{eq:dx-dmu-completo}
\frac{d x_t^\ast}{d\mu}
=
-\bigl(H_t^b\bigr)^{-1}\nabla_{x\mu}^2J_t^b(x_t^\ast,\mu).
\end{equation}
Taking the second coordinate yields
\begin{equation}
\label{eq:dgamma-dmu-completo}
\frac{d\gamma_t^\ast}{d\mu_t(\theta_{high})}
=
-e_\gamma'
\bigl(H_t^b\bigr)^{-1}
\nabla_{x\mu}^2J_t^b(x_t^\ast,\mu_t).
\end{equation}
By the net monotonicity condition~\eqref{eq:monotonia-neta-gamma}, the right-hand
side of~\eqref{eq:dgamma-dmu-completo} is strictly positive.

To interpret this condition, note that the coordinate of
$\nabla_{x\mu}^2J_t^b$ associated with $\gamma_t$ contains
\begin{equation}
\label{eq:efecto-neto-gamma}
\frac{\partial L_t}{\partial\gamma_t}
(b,\alpha_t^\ast,\gamma_t^\ast,\theta_{high},\iota_t)
-
\frac{\partial L_t}{\partial\gamma_t}
(b,\alpha_t^\ast,\gamma_t^\ast,\theta_{low},\iota_t)
+
\frac{\partial}{\partial\mu}
\left[
\frac{\partial \Pi_{t,b}^{S}}{\partial\gamma_t}
-\psi_H
\frac{\partial \Delta H_t}{\partial\gamma_t}
\right]_{x_t=x_t^\ast}.
\end{equation}
The first differential captures the material separation between types. Under
proportional risks,
\[
\frac{\tilde\lambda_j(t\mid\theta_{high})}
     {\tilde\lambda_j(t\mid\theta_{low})}
=
\exp\!\bigl(\beta_{K,j}(\theta_{high})
-\beta_{K,j}(\theta_{low})\bigr),
\]
so that greater structural distance in $j\in\{2,3\}$ increases the difference
in marginal losses associated with operational pressure. However, the deviation
premium and the informational subsidy also enter the first-order condition;
therefore the result is formulated as monotonicity of the net marginal effect
of the complete objective, not as a property of material loss alone. Under
~\eqref{eq:monotonia-neta-gamma}, increasing the posterior mass on the
high-risk type shifts the interior optimum toward a higher $\gamma_t^\ast$.
\end{proof}

% ============================================================================
\section{Equilibrium and implementability}
\label{sec:equilibrio}
% ============================================================================

\noindent This section introduces the dynamic mechanism with its PBE, the
notion of implementability with the constraints $IC^K$, $IR^K$, $IR^F$, and
proves the existence of the implementable equilibrium.

% ---------------------------------------------------------------------------

\begin{definition}[Dynamic mechanism and Perfect Bayesian Equilibrium]
\label{def:pbe-mechanism}
Consider a discrete-time Bayesian game with horizon
$\mathbb{T}=\{0,1,\ldots,\tau\}$, players $\mathcal{J}=\{S,K,F\}$, and captor
type $\theta_K\in\Theta_K$ according to common prior $\mu_0$ with
$\mu_0(\theta)>0$ for all $\theta$. A \textbf{dynamic mechanism} is described by:
\begin{itemize}
  \item \emph{Strategy profile:} $s=(s^S,s^K,s^F)$, where
  $s^S=\chi=(\chi_t)_{t\in\mathbb{T}}$ with
  $\chi_t:H_t\times\mathcal{R}
   \to\mathcal{A}^{S,\mathrm{disc}}\times\mathcal{A}^\alpha\times\mathcal{A}^\gamma$.
  \item \emph{Belief system:} $\mu=\{\mu_t\}_{t\ge 0}$,
  $\mu_t\in\Delta(\Theta_K)$.
  \item \emph{Ex-ante law:} $\ProbE^s$ induced by profile $s$, prior
  $\mu_0$, and the physical dynamics of captivity.
\end{itemize}

\noindent\textbf{Perfect Bayesian Equilibrium.} The pair $(s,\mu)$ is a
\textbf{PBE} if the following hold:
\begin{enumerate}
  \item[\textup{(a)}] \textbf{Sequential rationality.}
  At every node reached with positive probability, each agent $i$ solves
  its optimization program given $\mu_t$ and $s^{-i}$. For the state:
  \[
  \begin{aligned}
  (a_t^{S\ast},\alpha_t^\ast,\gamma_t^\ast)\in
  \arg\min_{(a_t^S,\alpha_t,\gamma_t)}
  \Bigl\{&
  \sum_\theta\mu_t(\theta)L_t(a_t^S,\alpha_t,\gamma_t,\theta,\iota_t)\\
  &+\Pi_{t,a^S}^S(\alpha_t,\gamma_t;\mu_t)
  -\psi_H\Delta H_t(\alpha_t,\gamma_t)\Bigr\}.
  \end{aligned}
  \]
  For $i\in\{K,F\}$:
  $a_t^{i\ast}\in\arg\max_{a_t^i}\mathcal{U}_t^i(a_t^i\mid\mathcal{I}_t^i,\chi,s^{-i})$.
  \item[\textup{(b)}] \textbf{On-path Bayesian consistency.}
  \[
  \mu_t(\theta)=\ProbE^s(\theta_K=\theta\mid\mathcal{F}_t^{\mathrm{pub}}),
  \qquad\forall\,\theta,\;t\ge 1.
  \]
\end{enumerate}
\end{definition}

\begin{remark}
Definition~\ref{def:pbe-mechanism} defines the PBE as an \emph{object}: the
pair $(s,\mu)$ satisfying~(a) and~(b). It does not assert existence nor that the
profile satisfies incentive and participation constraints; those properties are
the content of Theorem~\ref{thm:ebp-implementabilidad}.
\end{remark}

\bigskip

\begin{definition}[Nash-Bayes-$IC^K$,$IR^K$,$IR^F$ implementability]
\label{def:implementable-mechanism}
Let $(s^\ast,\mu^\ast)$ be a PBE (Definition~\ref{def:pbe-mechanism}). The
mechanism is \textbf{implementable Nash-Bayes-$IC^K$,$IR^K$,$IR^F$} if the
profile additionally satisfies the following conditions:

\smallskip
\noindent\textbf{Bayesian incentive compatibility of the captor ($IC^K$).}
For all $t\in\mathbb{T}$ and all rival mimicry $\theta_j\in\Theta_K$,
\[
\mathbb{E}_{\theta\sim\mu_t^\ast}
\!\left[
V^K_t\!\bigl(a^\ast(\theta)\mid\theta,\alpha_t^\ast,\gamma_t^\ast\bigr)
-V^K_t\!\bigl(a^\ast(\theta_j)\mid\theta,\alpha_t^\ast,\gamma_t^\ast\bigr)
\right]\ge 0.
\]
\noindent\textbf{Participation constraint of the captor ($IR^K$).}
\[
\mathbb{E}_{\theta\sim\mu_t^\ast}
\!\left[U^K_{\mathrm{rel}}(\theta)
-\max\bigl\{V^K_{\mathrm{cont},t}(\theta,\alpha_t^\ast,\gamma_t^\ast),
U^K_{\mathrm{kill}}(\theta,\theta_S)\bigr\}\right]\ge 0.
\]
\noindent\textbf{Participation constraint of the family ($IR^F$).}
\[
\mathcal{U}^F_t(a_{\mathrm{coop}})\ge\mathcal{U}^F_t(a_{\mathrm{col}}).
\]
The \textbf{feasible set of policies} is
\begin{equation}
\label{eq:gamma-factible}
\Gamma_t(\mu_t)
:=\Bigl\{(a_t^S,\alpha_t,\gamma_t)\in
\mathcal{A}^{S,\mathrm{disc}}\times[0,1]^2:
IC^K,\;IR^K,\;IR^F\text{ are satisfied}\Bigr\}.
\end{equation}
Its non-emptiness --- \textbf{robust feasibility} --- is the primitive
hypothesis of Theorem~\ref{thm:ebp-implementabilidad}.
\end{definition}

\begin{remark}
Definition~\ref{def:implementable-mechanism} states implementability as a
\emph{property} of the pair $(s^\ast,\mu^\ast)$: given the existence of the
PBE, one asks whether it additionally satisfies $IC^K$, $IR^K$, and $IR^F$. The
cycle criterion of \citet{Rochet1987} provides a useful algebraic
reformulation for auditing $IC^K$ when $\Theta_K$ is finite.
\end{remark}

\begin{remark}[Full support of the MDG and Bayesian operability]
The full support of $\ProbI{i}$ (hypothesis~(iii) of
Theorem~\ref{thm:ebp-implementabilidad}) guarantees strictly positive
probability at all reachable nodes, making Bayesian updating operative on the
path and reducing indeterminacy off the path.
\end{remark}

\bigskip

\begin{theorem}[Conditional implementability of the mechanism in PBE]
\label{thm:ebp-implementabilidad}
Consider the mechanism of Definition~\ref{def:pbe-mechanism} under the
following conditions:
\begin{itemize}
  \item[\textup{(i)}] $|\Theta_K|<\infty$ with strictly positive prior.
  \item[\textup{(ii)}] Finite action spaces $\mathcal{A}^i$ for all
  $i\in\{K,S,F\}$.
  \item[\textup{(iii)}] The law $\ProbI{i}$ of the MDG process assigns strictly
  positive probability to every action (\textbf{full support}).
  \item[\textup{(iv)}] $\mathcal{A}^\alpha=\mathcal{A}^\gamma=[0,1]$,
  compact and convex; $L_t$ and $\mathcal{U}_t^i$ are measurable, bounded, and
  continuous in the continuous decision variables. The physical outcome rule
  $\mathsf{G}_t$ is measurable, and the transition kernel $Q_t$ induced by
  $\mathsf{G}_t$ and by the competitive risk probabilities is continuous in
  those variables.
  \item[\textup{(v)}] \textbf{Robust feasibility.}
  For each $t$ and $\mu_t\in\Delta(\Theta_K)$, $\Gamma_t(\mu_t)$ is non-empty
  and compact; the sections $\Gamma_t^{a^S}(\mu_t)$ are convex; the
  correspondence $\mu_t\mapsto\Gamma_t(\mu_t)$ is continuous with compact
  values.
  \item[\textup{(vi)}] \textbf{Parametric separability.}
  \[
  V^K_t(a^\ast(\theta')\mid\theta_K,\alpha_t,\gamma_t)
  =
  \sum_jw_j(a^\ast(\theta'),\alpha_t,\gamma_t)
  \exp(\beta_{K,j}(\theta_K)\cdot z).
  \]
  \item[\textup{(vii)}] \textbf{Quasiconvexity of the restricted objective.}
  For each $t$, $\mu_t$, and $a_t^S$, the functional
  $(\alpha_t,\gamma_t)\mapsto\sum_\theta\mu_t(\theta)L_t+\Pi^S_{t,a^S}
  -\psi_H\Delta H_t$ is quasiconvex in $(\alpha_t,\gamma_t)$ over
  $\Gamma_t^{a^S}(\mu_t)$.
\end{itemize}
Then there exists $(s^\ast,\mu^\ast)$ such that:
\begin{enumerate}
  \item[\textup{(a)}] Constitutes a \textbf{PBE}
  (Definition~\ref{def:pbe-mechanism}).
  \item[\textup{(b)}] \textbf{Implements the social choice rule} $f_t$:
  $\chi_t^\ast(h_t)=(a_t^{S\ast},\alpha_t^\ast,\gamma_t^\ast)$ minimizes the
  expected social loss subject to $\Gamma_t(\mu_t^\ast)$.
  \item[\textup{(c)}] \textbf{Satisfies $IC^K$, $IR^K$, and $IR^F$} in profile
  $(s^\ast,\mu^\ast)$.
\end{enumerate}
\end{theorem}

\begin{proof}
The proof proceeds in five steps.

\paragraph{1. Existence of equilibrium.}
By~(ii), the behavioral strategy space $\Sigma^i$ of each agent is a finite
product of simplices, compact and convex. By~(v),
$\Gamma_t(\mu_t)$ is compact, non-empty, with convex sections and continuous
correspondence. The state's objective is continuous in $(\alpha_t,\gamma_t)$
by~(iv); by \textbf{Berge's maximum theorem} \citep{AliprantisBorder2006}, the
correspondence of optima $\Psi_t(\mu_t)$ is upper hemicontinuous and compact.
By~(vii), the functional is quasiconvex; the argmin of a quasiconvex function
over a compact convex set is convex, so $\Psi_t(\mu_t)$ has convex values. The
utilities of $K$ and $F$ are multilinear in behavioral strategies, so the best
response correspondence
\[
\mathrm{BR}(s)
:=\mathrm{BR}^S_\Gamma(s^{-S})\times\mathrm{BR}^K(s^{-K})
  \times\mathrm{BR}^F(s^{-F})
\]
has non-empty convex values. By \textbf{Kakutani's fixed-point theorem}
\citep{Kakutani1941}, there exists $s^\ast$ with
$s^\ast\in\mathrm{BR}(s^\ast)$.

\paragraph{2. On-path Bayesian consistency.}
By~(iii), $\ProbI{i}$ assigns positive probability to all actions compatible
with the physical dynamics. Bayes' formula
\[
\mu_t^\ast(\theta)
=\ProbE^{s^\ast}\!\bigl(\theta_K=\theta\mid\sigma(z,\zeta_0,\ldots,\zeta_{t-1})\bigr)
\]
is well defined at reached nodes and satisfies condition~(b) of
Definition~\ref{def:pbe-mechanism}.

\paragraph{3. Implementation of the social choice rule.}
At the fixed point, the state's strategy solves at each node on the path:
\[
(a_t^{S\ast},\alpha_t^\ast,\gamma_t^\ast)
\in\arg\min_{(a_t^S,\alpha_t,\gamma_t)\in\Gamma_t(\mu_t^\ast)}
\!\Bigl\{\sum_\theta\mu_t^\ast(\theta)L_t
+\Pi_{t,a^S}^S(\alpha_t,\gamma_t;\mu_t^\ast)
-\psi_H\Delta H_t(\alpha_t,\gamma_t)\Bigr\}.
\]
The objective is continuous by~(iv) and $\Gamma_t(\mu_t^\ast)$ is compact and
non-empty by~(v); by Weierstrass a solution exists.

\paragraph{4. Incentive compatibility ($IC^K$): cycle reformulation
\citep{Rochet1987}.}
$IC^K$ is imposed through the feasible set $\Gamma_t(\mu_t^\ast)$, not verified
ex post: equilibrium instruments belong to $\Gamma_t(\mu_t^\ast)$, which by~(v)
and Definition~\ref{def:implementable-mechanism} incorporates $IC^K$ as a
constitutive condition. For algebraic auditing with
$\Delta_{ij}^t:=V^K_t(a^\ast(\theta_j)\mid\theta_i)
-V^K_t(a^\ast(\theta_i)\mid\theta_i)$, the Bayesian version is equivalent to
$\mathbb{E}_{\theta_i\sim\mu_t^\ast}[\Delta_{ij}^t]\le 0$ for all $j$.
For every cycle $(\theta_{k_0},\ldots,\theta_{k_m},\theta_{k_0})$:
\[
\sum_{\ell=0}^{m}
\mu_t^\ast(\theta_{k_\ell})\,\Delta_{k_\ell,k_{\ell+1}}^t\le 0
\qquad(\text{indices mod }m+1).
\]

\paragraph{5. Participation constraints ($IR^K$, $IR^F$).}
By construction in Step~3, $(\alpha_t^\ast,\gamma_t^\ast)\in\Gamma_t(\mu_t^\ast)$,
a set that by~(v) simultaneously guarantees $IR^K$ and $IR^F$ in the
equilibrium profile.

Having completed the five steps, $(s^\ast,\mu^\ast)$ is a PBE that implements
$f_t$ and satisfies $IC^K$, $IR^K$, and $IR^F$ simultaneously.
\end{proof}

\begin{remark}[Separation of results: implementability and asymptotic identification]
Theorem~\ref{thm:ebp-implementabilidad} and Theorem
~\ref{thm:posterior-consistency} guarantee complementary properties. The first
guarantees existence of a PBE with implementable policy subject to
$\Gamma_t$: $IC^K$ enters through the definition of the feasible set. The
second guarantees that under that same equilibrium, if heterogeneity in
$\beta_{K,j}$ generates uniform separation in total variation, beliefs converge
to the true type almost surely in $\mathcal{T}$. The full support of the MDG
plays a role in both: it reduces the \emph{off-path} region and avoids
degeneracies of the likelihood.
\end{remark}

% ============================================================================
\section{Abstract structural identification (Kakutani criterion)}
\label{sec:id-abstracta}
% ============================================================================

\noindent This section isolates the measure-theoretic argument: under product
structure and recurrent separation in total variation, the criterion of
\citet{Kakutani1948} forces mutual singularity of trajectory laws.
Theorem~\ref{thm:posterior-consistency-appendix} states the same Bayesian
conclusion as Section~\ref{sec:id-asintotica}, replacing Lemma
~\ref{lem:hazard-separation} with the abstract hypotheses of the following
lemma.

% ---------------------------------------------------------------------------

\begin{lemma}[Product of experiments and singularity (Kakutani criterion)]
\label{lem:general-id-appendix}
Fix $\theta\neq\theta'$. Suppose that conditional on equilibrium policy and
on event $\mathcal{T}$, the sequence $\{Y_t\}_{t\ge 0}$ is independent and each
$Y_t$ takes values in a common finite support $\mathcal{Y}_t$. Let
$\rho_t:=\rho\bigl(p_\theta^{(t)},p_{\theta'}^{(t)}\bigr)$ be the
Bhattacharyya affinity. Then the affinity between the product laws on horizon
$T$ is
\begin{equation}
\label{eq:appendix-rho-product}
\rho\!\Bigl(\bigotimes_{t=0}^{T}p_\theta^{(t)},\,
            \bigotimes_{t=0}^{T}p_{\theta'}^{(t)}\Bigr)
=\prod_{t=0}^{T}\rho_t.
\end{equation}
If there exists an infinite set $\mathcal{S}_{\theta,\theta'}\subseteq\mathbb{N}$
with $\|p_\theta^{(t)}-p_{\theta'}^{(t)}\|_{\mathrm{TV}}\ge\underline\varepsilon>0$
for all $t\in\mathcal{S}_{\theta,\theta'}$, then
$\prod_{t\in\mathcal{S}}\rho_t=0$ and, by the criterion of
\citet{Kakutani1948}, the product measures over the infinite sequence are
\textbf{mutually singular} \citep{Shiryaev2019Probability}.
\end{lemma}

\begin{proof}
\textbf{Identity~\eqref{eq:appendix-rho-product}.}
Under independence,
\[
\rho\!\Bigl(\bigotimes_{t=0}^{T}p_\theta^{(t)},\,
            \bigotimes_{t=0}^{T}p_{\theta'}^{(t)}\Bigr)
=\sum_{(y_0,\ldots,y_T)}\prod_{t=0}^{T}\sqrt{p_\theta^{(t)}(y_t)p_{\theta'}^{(t)}(y_t)}
=\prod_{t=0}^{T}\sum_{y_t\in\mathcal{Y}_t}
  \sqrt{p_\theta^{(t)}(y_t)p_{\theta'}^{(t)}(y_t)}.
\]

\noindent\textbf{Singularity.}
The proof uses the standard inequality
$\|p-q\|_{\mathrm{TV}}\le\sqrt{2}\,\mathsf{H}(p,q)$
on finite spaces, with $\mathsf{H}^2(p,q)=1-\rho(p,q)$. If
$\|p-q\|_{\mathrm{TV}}\ge\underline\varepsilon$, then
\[
\rho(p,q)\le 1-\frac{\underline\varepsilon^{\,2}}{2}.
\]
At each $t\in\mathcal{S}_{\theta,\theta'}$,
$\rho_t\le 1-\underline\varepsilon^2/2<1$. As $\mathcal{S}$ is infinite, the
product contains infinitely many factors uniformly less than~$1$:
$\prod_{t\in\mathcal{S}}\rho_t=0$. By the criterion of \citet{Kakutani1948},
the infinite product measures are mutually singular
\citep{Shiryaev2019Probability}.
\end{proof}

\bigskip

\begin{theorem}[Posterior consistency under product-measure singularity]
\label{thm:posterior-consistency-appendix}
Suppose $\Theta_K$ finite, prior $\mu_0$ with full support, and PBE-consistent
Bayesian updating (Definition~\ref{def:beliefs-martingale}).
Let $\theta_K^\ast$ be the realized type and $\mathcal{T}$ the prolonged
captivity event (Definition~\ref{def:event-T-long}). Suppose that for all
$\theta\neq\theta_K^\ast$ the following hold:
\begin{enumerate}
  \item[\textup{(i)}] \textbf{Trajectory compatibility.}
  $\ProbC{\theta_K^\ast}(\mathcal{T})>0$; for each false type, either
  $\ProbC{\theta}(\mathcal{T})=0$, or signal laws are compared under the law
  restricted to $\mathcal{T}$.
  \item[\textup{(ii)}] \textbf{Product structure and recurrent TV separation.}
  The sequence $\{Y_t\}_{t\ge 0}$ satisfies product factorization and the
  existence of $\mathcal{S}_{\theta,\theta_K^\ast}$ and
  $\underline\varepsilon_\theta>0$ from Lemma~\ref{lem:general-id-appendix}.
  \item[\textup{(iii)}] \textbf{Signal-experiment alignment.}
  In $\mathcal{T}$, $\zeta_t$ contains $Y_t=(m_t,d_t)$, so that singularity of
  $\{Y_t\}$ transfers to $\{\zeta_t\}$ by monotonicity of singularity with
  respect to $\sigma$-algebras.
\end{enumerate}
Then, in $\mathcal{T}$,
\[
\mu_t(\theta)\xrightarrow{t\to\infty}0\;\;\forall\,\theta\neq\theta_K^\ast,
\qquad
\mu_t(\theta_K^\ast)\xrightarrow{t\to\infty}1,
\qquad\ProbC{\theta_K^\ast}\textup{-a.s.}
\]
Equivalently, $\{\mu_t\}$ converges weakly to $\delta_{\theta_K^\ast}$
$\ProbC{\theta_K^\ast}$-a.s.\ in $\mathcal{T}$.
\end{theorem}

\begin{proof}
Fix $\theta\neq\theta_K^\ast$.

\textbf{Case} $\ProbC{\theta}(\mathcal{T})=0$: the false type cannot generate
trajectories in $\mathcal{T}$; its posterior likelihood is annulled by survival
exclusion.

\textbf{Case} $\ProbC{\theta}(\mathcal{T})>0$: by~(ii), Lemma
~\ref{lem:general-id-appendix} implies that the infinite product laws are
mutually singular. By~(iii), that singularity transfers to the public flow
$\{\zeta_t\}$. There exists a tail event $A_\theta$ with
\[
\ProbC{\theta_K^\ast}(A_\theta\mid\mathcal{T})=1,\qquad
\ProbC{\theta}(A_\theta\mid\mathcal{T})=0.
\]
The likelihood ratio $R_t^\theta$ of the false type against the true type
satisfies $R_t^\theta\to 0$ $\ProbC{\theta_K^\ast}$-a.s.\ in $\mathcal{T}$
\citep{BlackwellDubins1962,Schervish1995Theory}. By Bayes,
\[
\mu_t(\theta)
\le\frac{\mu_0(\theta)}{\mu_0(\theta_K^\ast)}\,R_t^\theta\to 0
\qquad\ProbC{\theta_K^\ast}\text{-a.s.\ in }\mathcal{T}.
\]
As $\Theta_K$ is finite, all false coordinates converge to zero
simultaneously. Finally,
\[
1=\lim_{t\to\infty}\sum_\theta\mu_t(\theta)
=\lim_{t\to\infty}\mu_t(\theta_K^\ast)+0,
\]
so that $\mu_t(\theta_K^\ast)\to 1$ $\ProbC{\theta_K^\ast}$-a.s.\ in
$\mathcal{T}$.
\end{proof}

% ============================================================================
\section{Formal verification in Lean~4}
\label{sec:lean-formalization}
% ============================================================================

\begin{sloppypar}
\noindent
This replication package includes a standalone Lean~4 project in folder
\path{lean/} (main file \path{Appendix3Proofs.lean}) that
formalizes the logical dependency graph of this supplement. The development is
independent of earlier project-specific files and mirrors Sections
\ref{sec:main}--\ref{sec:id-abstracta}. Each declaration is tagged with its TeX
label where applicable. A successful build (\path{lake build}) yields no
\path{sorry} obligations. Build instructions appear in Section~\ref{sec:repro}.
\end{sloppypar}

\medskip
\noindent\textbf{Reading the status tags.}
Every Lean declaration carries one of two tags. \texttt{[Proved]} means the
statement is fully verified by the Lean~4 compiler against Mathlib: the proof
is machine-checked and admits no gaps. \texttt{[Hypothesis]} means the
statement is declared as an explicit interface and \emph{assumed}, not
re-proved, inside Lean; each such interface corresponds to a classical step
proved analytically in this supplement or in the cited literature. The
formalization therefore does \emph{not} replace the proofs in this document:
it certifies the finite-dimensional algebra and the logical assembly of the
identification argument, and makes every classical input explicit.

\medskip
\noindent\textbf{Proved in Lean (machine-checked).}
\begin{itemize}
  \item \textbf{Section~\ref{sec:main}.}
  Lemma~\ref{lem:I0-closed} and Corollary~\ref{cor:trend_equivalence}.
  \item \textbf{Section~\ref{sec:info-bayes}.}
  Public filtration, predictive likelihood, and martingale beliefs;
  Lemma~\ref{lem:asymptotic-beliefs} (bounded martingale with a.s.\ limit in
  $[0,1]$, via Doob's theorem in Mathlib); and the normalization closure of
  step~3 of Theorem~\ref{thm:posterior-consistency}: false-type collapse
  forces $\mu_t(\theta^\ast)\to 1$.
  \item \textbf{Section~\ref{sec:mdg}.}
  Continuation probability of Definition~\ref{def:desenlace-fisico-mt}; Cox
  hazard-ratio algebra (distinct $\beta_{K,j}$ $\Rightarrow$ ratio
  $\neq 1$); the predictive masses $p_\theta^{(t)}$ form probability masses;
  local TV separation --- Lemma~\ref{lem:hazard-separation}(i) --- for those
  masses under $\beta_{K,j}$ separation and a non-knife-edge total hazard;
  logit axioms (A1)~full support and (A2)~plan dominance of
  Remark~\ref{rem:axiomas-logit}.
  \item \textbf{Section~\ref{sec:riesgos}.}
  Sign certificate of Proposition~\ref{prop:expected-captivity-gamma}:
  negativity of~\eqref{eq:dE-dgamma} and the sign identity under
  $\kappa_h\gtrless 0$.
  \item \textbf{Section~\ref{sec:id-asintotica}.}
  Logical assembly of Lemma~\ref{lem:hazard-separation}: conditions (a)/(b)
  imply Definition~\ref{def:asymptotic-id}; packaging of recurrent uniform TV
  margins; and the Bayes factorization of step~2 of
  Theorem~\ref{thm:posterior-consistency}: the bound
  $\mu_t(\theta)\le(\mu_0(\theta)/\mu_0(\theta^\ast))\,R_t^\theta$ forces
  $\mu_t(\theta)\to 0$ whenever $R_t^\theta\to 0$.
  \item \textbf{Section~\ref{sec:id-abstracta}.}
  Bhattacharyya product identity~\eqref{eq:appendix-rho-product}; the
  Hellinger--TV bound $\rho\le 1-\underline\varepsilon^2/2$; and vanishing of
  the product affinity under recurrent separation, i.e., the entire algebraic
  part of Lemma~\ref{lem:general-id-appendix}.
  \item \textbf{Identification bridge.}
  Cox predictive masses plus a recurrent TV margin imply asymptotic
  identification, conditional only on the two Kakutani interfaces listed
  below.
\end{itemize}

\medskip
\noindent\textbf{Not proved in Lean (explicit hypotheses).}
The following classical steps are stated as interfaces and assumed inside
Lean, mainly because infinite product measures and their convergence theory
are not available in Mathlib. Each one is proved analytically where indicated.
\begin{itemize}
  \item \emph{Kakutani mutual singularity} --- vanishing product affinity
  implies mutually singular trajectory laws: proof of
  Lemma~\ref{lem:general-id-appendix} and \citet{Kakutani1948}.
  \item \emph{Singularity transfer} $Y_t\Rightarrow\zeta_t$: item~(iii) of
  Theorem~\ref{thm:posterior-consistency-appendix}.
  \item \emph{Singularity implies vanishing likelihood ratio}: step~2 of
  Theorem~\ref{thm:posterior-consistency}
  \citep{BlackwellDubins1962,Schervish1995Theory}.
  \item \emph{Survival exclusion}, $\ProbC{\theta}(\mathcal{T})=0$ implies
  collapse: step~2 of Theorem~\ref{thm:posterior-consistency}.
  \item \emph{Maximum entropy axiom (A3)} of Remark~\ref{rem:axiomas-logit}:
  \citet{Jaynes1957Information}.
  \item \emph{Implicit-function step} of
  Proposition~\ref{prop:optimal-gamma-beliefs}: proof in
  Section~\ref{sec:politica}.
  \item \emph{Existence of the implementable PBE},
  Theorem~\ref{thm:ebp-implementabilidad}: fixed-point proof in
  Section~\ref{sec:equilibrio} \citep{Kakutani1941}.
\end{itemize}
Table~\ref{tab:lean-status} maps these hypotheses to their Lean identifiers.

\begin{table}[htbp]
\centering
\scriptsize
\caption{Lean formalization status for selected supplement objects}
\label{tab:lean-status}
\setlength{\tabcolsep}{3pt}
\renewcommand{\arraystretch}{1.15}
\begin{tabular}{@{}%
  >{\raggedright\arraybackslash}p{0.37\linewidth}%
  >{\raggedright\arraybackslash}p{0.41\linewidth}%
  >{\raggedright\arraybackslash}p{0.12\linewidth}@{}}
\toprule
TeX object & Lean identifier & Status \\
\midrule
Lemma~\ref{lem:general-id-appendix} (algebraic part)
  & \texttt{generalIdAppendix\-Lemma} & Proved \\
Kakutani (1948) mutual singularity
  & \texttt{kakutani\-Mutual\-Singularity} & Hyp. \\
Outcome-trajectory laws on $\mathcal{T}$
  & \texttt{kakutani\-Signal\-Singularity} & Hyp. \\
Singularity transfer $Y_t\Rightarrow\zeta_t$
  & \texttt{signal\-Experiment\-Alignment} & Hyp. \\
Singularity $\Rightarrow$ vanishing likelihood ratio
  & \texttt{singularity\-Implies\-Likelihood\-Vanishing} & Hyp. \\
$\ProbC{\theta}(\mathcal{T})=0$ $\Rightarrow$ collapse
  & \texttt{zeroMass\-Posterior\-Collapse} & Hyp. \\
Remark~\ref{rem:axiomas-logit}, axiom (A3)
  & \texttt{logitProb\-\_max\-Entropy\-\_hyp} & Hyp. \\
Proposition~\ref{prop:optimal-gamma-beliefs}
  & \texttt{prop\-\_optimal\-\_gamma\-\_beliefs} & Hyp. \\
Theorem~\ref{thm:ebp-implementabilidad}
  & \texttt{conditional\-Implementability} & Hyp. \\
\bottomrule
\end{tabular}
\tabnote\ \texttt{[Proved]}: machine-checked in Lean~4. \texttt{[Hyp.]}:
declared interface; proved analytically in this supplement or cited literature.
\tabsource
\end{table}

% ============================================================================
\section{Reproducibility}
\label{sec:repro}
% ============================================================================

\begin{sloppypar}
All analytical proofs in this supplement are self-contained in
\path{tex/Appendix_3.tex}. Machine-checked algebraic steps are reproduced by
the Lean~4 project in \path{lean/}. Set the working directory to
\texttt{Appendix\_3/} and run \path{replicate.sh}. See \path{README.txt} at the
package root for step-by-step instructions. Requirements: Lean~4 (toolchain in
\path{lean-toolchain}) and Lake; Mathlib is fetched automatically on first
build (network required; estimated first-build time 15--45 minutes).
\end{sloppypar}

\begin{enumerate}
\item \path{replicate.sh} at the package root runs \path{lake build} in
  \path{lean/} and writes a log to \path{lean/build.log}.
\item Main formal file: \path{lean/Appendix3Proofs.lean}. Project metadata:
  \path{lakefile.lean}, \path{lean-toolchain}, \path{lake-manifest.json}.
\item Section~\ref{sec:lean-formalization} and Table~\ref{tab:lean-status}
  document which results are \texttt{[Proved]} vs.\ \texttt{[Hyp.]}.
\item (Optional) Compile \path{tex/Appendix_3.pdf} from \path{tex/} with
  \texttt{pdflatex} and \texttt{bibtex}.
\end{enumerate}

\begin{sloppypar}
Section~\ref{sec:lean-formalization} (Table~\ref{tab:lean-status}) documents
which results are machine-checked vs.\ hypothetical. All paths are relative to
\texttt{Appendix\_3/}; see \path{README.txt} for the full build workflow.
\end{sloppypar}

\bibliographystyle{econsoc}
\bibliography{references}

\end{document}
